%%%%%%%%%%%%%%%%%%%%%%%%%%%%%%%%%%%%%%%%%
% Fancyslides Presentation
% LaTeX Template
% Version 1.0 (30/6/13)
%
% This template has been downloaded from:
% http://www.LaTeXTemplates.com
%
% The Fancyslides class was created by:
% Paweł Łupkowski (pawel.lupkowski@gmail.com)
%
% License:
% CC BY-NC-SA 3.0 (http://creativecommons.org/licenses/by-nc-sa/3.0/)
%
%%%%%%%%%%%%%%%%%%%%%%%%%%%%%%%%%%%%%%%%%

%----------------------------------------------------------------------------------------
%	PACKAGES AND OTHER DOCUMENT CONFIGURATIONS
%----------------------------------------------------------------------------------------

\documentclass{fancyslides}

\usepackage[utf8]{inputenc} % Allows the usage of non-english characters
\usepackage{librebaskerville} % Use the Times font
\usepackage{booktabs} % Allows the use of \toprule, \midrule and \bottomrule in tables
\usepackage[percent]{overpic}
\graphicspath{{images/}} % Location of the slide background and figure files

% Beamer options - do not change
\usetheme{default} 
\setbeamertemplate{navigation symbols}{} % Disable the slide navigation buttons on the bottom of each slide
\setbeamercolor{structure}{fg=\yourowntexcol} % Define the color of titles and fixed text elements (e.g. bullet points)
\setbeamercolor{normal text}{fg=\yourowntexcol} % Define the color of text in the presentation

%------------------------------------------------
% COLORS
% The following colors are predefined in this class: white, black, gray, blue, green and orange

% Define your own color as follows:
\definecolor{cdv}{HTML}{3FBC7C}

\newcommand{\structureopacity}{0.75} % Opacity (transparency) for the structure elements (boxes and circles)

\newcommand{\strcolor}{cdv} % Set the color of structure elements (boxes and circles)
\newcommand{\yourowntexcol}{white} % Set the text color

\begin{document}

%------------------------------------------------

\fbckg{Kacheln.jpg} % Slide background image
\begin{frame}
\begin{center}
\begin{overpic}[scale=0.8]{CDV_Logo.eps}
\put (18,10) {Ein Vortrag von Felix Potthast und Daniel Busch}
\end{overpic}
\end{center}
\end{frame}

%------------------------------------------------

\fbckg{Kacheln.jpg} % Slide background image
\begin{frame}
\misc{
\Huge
Inhalt\\
\Large
\begin{flushleft}
\fitem{Was ist das Coding da Vinci?}
\fitem{Geschichte des Coding da Vinci}
\fitem{Beispielprojekte}
\fitem{Unser Projekt: Pyanola}
\fitem{Ausblick}
\fitem{Links}
\end{flushleft}
}
\end{frame}

%------------------------------------------------

\fbckg{Kacheln.jpg} % Slide background image
\begin{frame}
\misc{
\Huge
Was ist das Coding da Vinci?\\
\Large
\begin{flushleft}
\fitem{deutsche Kulturinstitutionen stellen offene Daten bereit}
\fitem{Kultur-Hackathon}
\fitem{Auftaktveranstaltung, dann 10 Wochen Sprint}
\fitem{Preisverleihung in Berlin}
\fitem{Ergebnisse stehen unter freier Lizenz}
\end{flushleft}
}
\end{frame}

%------------------------------------------------

\fbckg{Kacheln.jpg} % Slide background image
\begin{frame}
\misc{
\Huge
Veranstalter\\
\begin{center}
\includegraphics[scale=0.1]{okfde.eps}\, 
\includegraphics[scale=0.2]{WMD.eps}\,  \\
\smallskip
\includegraphics[scale=0.2]{DDB.eps}\, 
\includegraphics[scale=0.3]{digis.png}\, 
\end{center}
}
\end{frame}

%------------------------------------------------

\fbckg{Kacheln.jpg} % Slide background image
\begin{frame}
\misc{
\Huge
Geschichte des Coding da Vinci\\
\Large
\begin{flushleft}
\fitem{Erstmals 2014}
\fitem{Damals 16 Kulturinstitutionen mit 24 Datensets}
\fitem{17 Projekte fertig gestellt}
\fitem{Wiederholung 2015}
\fitem{Diesesmal 33 Kulturinstitutionen mit 47 Datensets}
\fitem{20 Projekte fertig gestellt}
\end{flushleft}
}
\end{frame}

%------------------------------------------------

\fbckg{Kacheln.jpg} % Slide background image
\begin{frame}
\misc{
\Huge
Beispielprojekte\\
\Large
\begin{flushleft}
\fitem{Cyberbeetle - Käfer Roboter}
\fitem{inside - 19xx - Website über verbotene Bücher}
\fitem{imperii-viz - Website über Kaiser und Könige}
\fitem{floradex - App zur Bestimmung von Pflanzen}
\fitem{Bilder der Revolte - Fotos der 68er Bewegung}
\end{flushleft}
}
\end{frame}

%------------------------------------------------

\fbckg{Kacheln.jpg} % Slide background image
\begin{frame}
\misc{
\Huge
Unser Projekt: Pyanola\\
\Large
\begin{flushleft}
\fitem{Klavierrollen vom deutschen Museum}
\fitem{Analyse, Generierung von MIDI Dateien}
\fitem{Präsentation auf einer Website}
\fitem{Status: fast fertig ;)}
\end{flushleft}
}
\end{frame}

%------------------------------------------------

\fbckg{Kacheln.jpg} % Slide background image
\begin{frame}
\misc{
\Huge
Ausblick\\
\Large
\begin{flushleft}
\fitem{2016 Kein Coding Da Vinci}
\fitem{2017 "in neuem Gewand"}
\end{flushleft}
}
\end{frame}

%------------------------------------------------

\fbckg{Kacheln.jpg} % Slide background image
\begin{frame}
\pointedsl{Fragen?}
\end{frame}

%------------------------------------------------

\fbckg{Kacheln.jpg} % Slide background image
\begin{frame}
\misc{
\Huge
Links\\
\Large
\begin{flushleft}
\fitem{http://codingdavinci.de/}
\fitem{http://bilder-der-revolte.de/}
\fitem{http://codingdavinci.de/projekte-2014/?project\_id=1}
\fitem{http://wesuku.de/floradex/www/}
\fitem{http://imperii-viz.de/}
\end{flushleft}
}
\end{frame}

%------------------------------------------------

\end{document}
